\documentclass[aspectratio=169]{beamer}

% ---------- ENCODING & FONTS ----------
\usepackage[utf8]{inputenc}
\usepackage[T1]{fontenc}

% ---------- THEME & LOOK ----------
\usetheme{Madrid}
\usecolortheme{seahorse}
\usefonttheme{professionalfonts}
\setbeamertemplate{navigation symbols}{}
\setbeamercovered{transparent}
\setbeamertemplate{itemize items}[circle]
\setbeamertemplate{enumerate items}[default]
\setlength{\parskip}{0.25em}

% ---------- COLORS ----------
\usepackage{xcolor}
\definecolor{accent}{RGB}{0,92,184}
\definecolor{soft}{RGB}{233,244,255}
\setbeamercolor{structure}{fg=accent}
\setbeamercolor{title}{fg=white,bg=accent}
\setbeamercolor{frametitle}{fg=accent}
\setbeamercolor{block title}{bg=accent!12,fg=accent}
\setbeamercolor{block body}{bg=soft}

% ---------- PACKAGES ----------
\usepackage{booktabs}
\usepackage{amsmath, amssymb}
\usepackage{array}
\usepackage{hyperref}
\hypersetup{colorlinks=true,linkcolor=accent,urlcolor=accent}

% ---------- SAFE TRANSITION FALLBACK ----------
\makeatletter
\@ifundefined{transfade}{\newcommand{\transfade}[1][]{}}{}
\makeatother

% ---------- PLACEHOLDER BOX (FIGURES/DIAGRAMS) ----------
\newcommand{\placeholderbox}[2][0.82\linewidth]{%
  \begin{center}
    \fbox{\begin{minipage}[c][0.46\textheight][c]{#1}
      \centering \vspace{0.25em}
      \textbf{#2}\\[0.6em]
      \emph{Insert figure/diagram here later.}
    \end{minipage}}
  \end{center}
}

% ---------- TITLE ----------
\title{\textbf{Economics: Foundations and Models  \vspace{.5cm}\\  }}
\subtitle{ {\large ECON 101 - Macroeconomics}}
%\institute{UNBC}
\author{Muhebullah Karimzada}
\date{Winter 2026}


\begin{document}

% ============================================================
% TITLE
% ============================================================
\begin{frame}[plain]
  \titlepage
\end{frame}
%
 %============================================================
\begin{frame}{Outline}
\transfade
\begin{itemize}
  \item \textbf{1.1} Three Key Economic Ideas
  \item \textbf{1.2} The Economic Problems All Societies Must Solve
  \item \textbf{1.3} Economic Models
  \item \textbf{1.4} Microeconomics and Macroeconomics
  \item \textbf{1.5} The Language of Economics
\end{itemize}
\end{frame}

% ============================================================
\section{What Is Economics?}
% ============================================================

\begin{frame}{What Is Economics? (1 of 2)}
\transfade
\begin{itemize}
  \item People make \textbf{choices} as they try to attain their goals.
  \item Choices are necessary because we live in a world of \alert{scarcity}.
  \item \textbf{Scarcity:} a situation in which unlimited wants exceed the limited resources available to fulfill those wants.
  \item \textbf{Economics} is the study of the choices people make to attain their goals, given their scarce resources.
\end{itemize}
\end{frame}

\begin{frame}{What Is Economics? (2 of 2)}
\transfade
\begin{itemize}
  \item Economists study these choices using \textbf{economic models}:
  \begin{itemize}
    \item \textbf{Economic model:} a simplified version of reality used to analyze real-world economic situations.
  \end{itemize}
  \item Typical economics questions:
  \begin{itemize}
    \item How are the prices of goods and services determined?
    \item How do government spending and taxation affect the economy?
    \item Why do governments control the prices of some goods?
    \item Why are some countries wealthier than others?
  \end{itemize}
\end{itemize}
\end{frame}

\begin{frame}{Everyday Examples of Scarcity}
\transfade
\begin{itemize}
  \item \textbf{Your time:} cannot work, study, and relax at the same time.
  \item \textbf{Government budgets:} more spending on health care may mean less on education or infrastructure.
  \item \textbf{Environmental trade-offs:} more production today may mean less clean air tomorrow.
\end{itemize}
\end{frame}



% ============================================================
\section{Three Key Economic Ideas (1.1)}
% ============================================================

\begin{frame}{Three Key Economic Ideas}
\transfade
\begin{itemize}
  \item People are \textbf{rational}.
  \item People respond to \textbf{incentives}.
  \item Optimal decisions are made at the \textbf{margin}.
\end{itemize}
\end{frame}

\begin{frame}{1. People Are Rational}
\transfade
\begin{itemize}
  \item Economists generally assume that people are \textbf{rational}.
  \item \textbf{Rational:} people make choices based on what they believe will make them happy.
  \item They do not deliberately do things to make themselves worse off.
  \item Rational decision makers weigh \alert{costs} and \alert{benefits}.
\end{itemize}
\end{frame}

\begin{frame}{Example: A \$5 Cup of Coffee}
\transfade
\begin{itemize}
  \item You will not buy a \$5 cup of coffee unless you believe the benefit is worth at least:
  \begin{itemize}
    \item the time and effort it takes you to earn \$5, \textbf{plus}
    \item the value of any other item you could have bought with the \$5.
  \end{itemize}
  \item \textbf{Hubbard-style interpretation:} Think in terms of \alert{opportunity cost}:
  \begin{itemize}
    \item If the best alternative use of \$5 is a bus pass or saving for tuition, you compare happiness from coffee vs.\ that alternative.
  \end{itemize}
\end{itemize}
\end{frame}

\begin{frame}{2. People Respond to Incentives}
\transfade
\begin{itemize}
  \item As \textbf{costs} and \textbf{benefits} change, so do people's actions.
  \item Example: Do Canadian banks choose to get robbed?
  \begin{itemize}
    \item Bullet-resistant shields cost \$20{,}000.
    \item Typical loss due to an armed robbery is \$1{,}200.
    \item Armed robbery is rare.
    \item \textbf{Result:} Many banks choose not to install bullet-resistant shields.
  \end{itemize}
\end{itemize}
\end{frame}

\begin{frame}{More Incentive Examples}
\transfade
\begin{itemize}
  \item \textbf{Students:} Higher course weight or harder marking may increase study time.
  \item \textbf{Firms:} A carbon tax raises the cost of emissions, encouraging cleaner technologies.
  \item \textbf{Households:} Higher interest rates increase the reward to saving but raise borrowing costs.
\end{itemize}
\end{frame}

\begin{frame}{3. Optimal Decisions Are Made at the Margin}
\transfade
\begin{itemize}
  \item Most decisions involve doing a \textbf{little more} or a \textbf{little less} of something.
  \item Example: Should you watch an extra hour of TV or study instead?
  \item \textbf{Marginal cost (MC):} the additional cost from a small increase in activity.
  \item \textbf{Marginal benefit (MB):} the additional benefit from that small increase.
  \item Comparing MC and MB is called \textbf{marginal analysis}.
\end{itemize}
\end{frame}

\begin{frame}{Worked Example: Study vs.\ TV}
\transfade
\small
Suppose:
\begin{itemize}
  \item Extra hour of TV: MB = 8 ``happiness points'', MC = 0 points.
  \item Extra hour of study: increases exam grade by 2 percentage points, which you value at 12 points of happiness (less stress, better GPA).
  \item \textbf{Decision:} If you switch from TV to study:
  \begin{itemize}
    \item MB (study) = 12, MC (giving up TV) = 8.
    \item Net gain = 4 points \(\Rightarrow\) better to study the extra hour.
  \end{itemize}
\end{itemize}
\end{frame}

% ============================================================
\section{Economic Problems All Societies Must Solve (1.2)}
% ============================================================

\begin{frame}{Three Fundamental Questions}
\transfade
\textbf{1.2 Learning Objective}\\[0.3em]
Discuss how an economy answers these questions:
\begin{itemize}
  \item \textbf{What} goods and services will be produced?
  \item \textbf{How} will the goods and services be produced?
  \item \textbf{Who} will receive the goods and services produced?
\end{itemize}
\end{frame}

\begin{frame}{1. What Goods and Services Will Be Produced?}
\transfade
\begin{itemize}
  \item All societies must decide which goods and services to produce.
  \item Producing more of one thing means producing less of another.
  \item This is a \textbf{trade-off}, resulting from the scarcity of productive resources.
  \item \textbf{Opportunity cost:} the value of the best alternative given up.
  \item Example: The opportunity cost of increased funding for health care might be fewer university scholarships.
\end{itemize}
\end{frame}



\begin{frame}{2. How Will the Goods Be Produced?}
\transfade
\begin{itemize}
  \item A firm may have multiple methods for producing the same output.
  \item \textbf{Example 1 (Music production):}
  \begin{itemize}
    \item Hire a great singer and use standard production techniques; or
    \item Hire a mediocre singer and use Auto-Tune to correct inaccuracies.
  \end{itemize}
  \item \textbf{Example 2 (Manufacturing labour):}
  \begin{itemize}
    \item As labour costs rise, firms may substitute capital (machines) for labour; or
    \item Move production to a country with cheaper labour.
  \end{itemize}
\end{itemize}
\end{frame}

\begin{frame}{3. Who Will Receive the Goods and Services?}
\transfade
\begin{itemize}
  \item Who gets what is closely tied to \alert{income distribution}.
  \item People with higher incomes can buy more goods and services.
  \item Governments can change the income distribution through:
  \begin{itemize}
    \item \textbf{Taxes}
    \item \textbf{Transfers} (e.g., EI, child benefits, social assistance)
  \end{itemize}
\end{itemize}
\end{frame}

\begin{frame}{Centrally Planned vs.\ Market vs.\ Mixed Economies}
\transfade
\begin{itemize}
  \item \textbf{Centrally planned economy:} government decides how economic resources will be allocated.
  \item \textbf{Market economy:} decisions of households and firms interacting in markets allocate resources.
  \item \textbf{Mixed economy:} most decisions are market-based, but government plays a significant role.
  \item Canada today is a \textbf{mixed economy}:
  \begin{itemize}
    \item Hospitals: centrally planned by provincial governments.
    \item Fast food: organized by markets.
  \end{itemize}
\end{itemize}
\end{frame}

\begin{frame}{Efficiency and Equity}
\transfade
\begin{itemize}
  \item Market economies tend to be more \textbf{efficient} than centrally planned economies.
  \item \textbf{Productive efficiency:} goods and services produced at the lowest possible cost.
  \item \textbf{Allocative efficiency:} production in accordance with consumer preferences:
  \begin{itemize}
    \item Each good produced up to the point where last unit's \textbf{MB} equals its \textbf{MC}.
  \end{itemize}
  \item \textbf{Equity:} the fair distribution of economic benefits.
\end{itemize}
\end{frame}

\begin{frame}{Source of Efficiency}
\transfade
\begin{itemize}
  \item \textbf{Productive efficiency} comes from \alert{competition} among firms.
  \item \textbf{Allocative efficiency} arises from \alert{voluntary exchange}:
  \begin{itemize}
    \item When buyer and seller both benefit from trade.
    \item Transactions continue until no further mutual gains are possible.
  \end{itemize}
\end{itemize}
\end{frame}

\begin{frame}{When Markets May Fail}
\transfade
\begin{itemize}
  \item Markets may not result in fully efficient outcomes:
  \begin{itemize}
    \item It takes time to discover the most efficient production methods.
    \item Governments may interfere with market outcomes (taxes, controls).
    \item Market outcomes can ignore those not directly involved:
    \begin{itemize}
      \item \textbf{Externalities} (e.g., pollution).
    \end{itemize}
  \end{itemize}
  \item Economic efficiency is not the only objective; \textbf{equity} matters too.
\end{itemize}
\end{frame}

% ============================================================
\section{Economic Models (1.3)}
% ============================================================

\begin{frame}{Economic Models}
\transfade
\begin{itemize}
  \item Economists develop models to analyze real-world issues.
  \item Building a model often follows these steps:
  \begin{itemize}
    \item Decide on the assumptions.
    \item Formulate a testable hypothesis.
    \item Use economic data to test the hypothesis.
    \item Revise the model if it fails to explain the data well.
    \item Retain the revised model for similar questions in the future.
  \end{itemize}
\end{itemize}
\end{frame}

\begin{frame}{Important Features of Economic Models}
\transfade
\begin{itemize}
  \item \textbf{Assumptions:} every model needs them to be tractable and useful.
  \item \textbf{Hypothesis testing:} good models generate testable predictions.
  \item \textbf{Economic variable:} something measurable that can take different values (e.g., income, unemployment rate).
\end{itemize}
\end{frame}

\begin{frame}{The Scientific Nature of Economics}
\transfade
\begin{itemize}
  \item Economists try to mimic the scientific method:
  \begin{itemize}
    \item Develop models
    \item Derive predictions
    \item Compare with data
  \end{itemize}
  \item Economics is a \textbf{social science}:
  \begin{itemize}
    \item People may have strong views about what results ``should'' be.
  \end{itemize}
\end{itemize}
\end{frame}

\begin{frame}{Positive vs.\ Normative Analysis}
\transfade
\begin{itemize}
  \item \textbf{Positive analysis:} based on facts and logic.
  \begin{itemize}
    \item ``If the minimum wage rises, some low-skill employment will fall.''
  \end{itemize}
  \item \textbf{Normative analysis:} based on value judgments.
  \begin{itemize}
    \item ``The minimum wage should be higher to ensure a living standard.''
    \item Words like ``should'' or ``ought'' often signal normative statements.
  \end{itemize}
  \item Economics \textbf{tries} to focus on positive analysis, while recognizing normative debates.
\end{itemize}
\end{frame}

\begin{frame}{Apply the Concept: Cap on House Prices (1 of 2)}
\transfade
\begin{itemize}
  \item Consider a government \textbf{cap} on house prices.
  \item \textbf{Potential benefits:}
  \begin{itemize}
    \item Lower home prices for buyers.
    \item Lower housing costs free up income for other spending.
  \end{itemize}
  \item \textbf{Potential costs:}
  \begin{itemize}
    \item Fewer houses built.
    \item Fewer jobs in construction and related sectors.
  \end{itemize}
\end{itemize}
\end{frame}

\begin{frame}{Apply the Concept: Cap on House Prices (2 of 2)}
\transfade
\begin{itemize}
  \item Both benefits and costs are \textbf{positive} effects we can try to measure.
  \item How much each effect matters is a \textbf{normative} judgment.
  \item This example illustrates:
  \begin{itemize}
    \item Trade-offs and opportunity cost.
    \item Distinction between positive and normative analysis.
  \end{itemize}
\end{itemize}
\end{frame}

% ============================================================
\section{Microeconomics and Macroeconomics (1.4)}
% ============================================================

\begin{frame}{Microeconomics and Macroeconomics}
\transfade
\textbf{1.4 Learning Objective}\\[0.3em]
Distinguish between microeconomics and macroeconomics.
\begin{itemize}
  \item \textbf{Microeconomics:} the study of how households and firms make choices, how they interact in markets, and how governments attempt to influence their choices.
  \item \textbf{Macroeconomics:} the study of the economy as a whole, including inflation, unemployment, and economic growth.
\end{itemize}
\end{frame}

\begin{frame}{Examples: Micro vs.\ Macro Questions}
\transfade
\begin{itemize}
  \item \textbf{Micro:}
  \begin{itemize}
    \item Why did Netflix raise its subscription price?
    \item How will a tax on sugary drinks affect soft drink sales?
  \end{itemize}
  \item \textbf{Macro:}
  \begin{itemize}
    \item Why does Canada experience recessions?
    \item What causes inflation to rise or fall?
  \end{itemize}
\end{itemize}
\end{frame}



% ============================================================
\section{The Language of Economics (1.5)}
% ============================================================

\begin{frame}{The Language of Economics (1 of 2)}
\transfade
\textbf{1.5 Learning Objective}\\[0.3em]
Define important economic terms.
\begin{itemize}
  \item Economics has its own terms and jargon with specific, precise meanings.
  \item Sometimes these terms are used differently than in everyday language.
\end{itemize}
\end{frame}

\begin{frame}{The Language of Economics (2 of 2)}
\transfade
\begin{itemize}
  \item \textbf{Technology:} the processes a firm uses to turn inputs into outputs of goods and services.
  \item \textbf{Capital:} durable manufactured goods used to produce other goods and services.
  \item Pay close attention to terms defined in class and in the textbook:
  \begin{itemize}
    \item Using terms correctly helps you think like an economist.
    \item On exams, definitions are often tested directly or indirectly.
  \end{itemize}
\end{itemize}
\end{frame}

% ============================================================
\section{Summary}
% ============================================================

\begin{frame}{Summary}
\transfade
\begin{itemize}
  \item Economics is about how individuals and societies make choices under \textbf{scarcity}.
  \item Three key ideas:
  \begin{itemize}
    \item People are rational.
    \item People respond to incentives.
    \item Optimal decisions are made at the margin.
  \end{itemize}
  \item Societies face three fundamental questions: what, how, and who.
  \item Market, centrally planned, and mixed economies differ in how they answer these questions.
\end{itemize}
\end{frame}

\begin{frame}{Summary (continued)}
\transfade
\begin{itemize}
  \item Economic models simplify reality to generate testable predictions.
  \item Economists distinguish between \textbf{positive} (what is) and \textbf{normative} (what ought to be) analysis.
  \item Microeconomics focuses on individual markets; macroeconomics looks at the economy as a whole.
  \item Mastering the \textbf{language of economics} (terms like scarcity, opportunity cost, technology, capital) is essential for success in this course.
\end{itemize}
\end{frame}

% ============================================================
\begin{frame}
\centering
\Huge \textbf{Thank you!}
\end{frame}

\end{document}
